% ============================================================
% DRAFT NOTE (method)
% Contains: The forecaster's causal inputs, the rate-based multi-horizon
% architecture, and the comparators.
% Written now: The selected model per the frozen current-state-damage-v1
% confirmation (2026-09-01): four XGBoost segment-rate regressors.
% Pending: Exact hyperparameters and feature-inventory hash from
% results/recovered/current-state-damage-v1/protocol_lock.json; precise
% definition of the action-clock comparator (verify against the lock).
% Sources: current-state-damage-v1 result; pivot discussion
% ============================================================
\section{HERALD: A Causal Multi-Horizon Forecaster}

\todo{STALE SECTION: written for the divergence objective; to be rewritten for the quality-risk target. Do not review yet.}
\label{sec:method}

\paragraph{Inputs.}
At step $t$ the forecaster sees only causally available, zero-cost signals,
in four groups. First, the active compression action: the compressor
identity, the removal ratio, and their combination, which set the baseline
risk of the current configuration. Second, generation position, as
$\log(1 + t)$, since compression behavior can change as the generated cache
grows. Third, the current state of the compressed model's next-token
distribution: entropy, top-1 and top-5 probability, effective vocabulary
size, alternative-token uncertainty, average log-probability, tail
probability mass, and logit range. Fourth, recent causal dynamics of that
distribution: entropy change, step-to-step KL divergence, top-10 overlap
between successive steps, and, for six principal sensors, trailing means and
standard deviations over 8- and 32-token windows and exponential moving
averages at half-lives 8 and 32, all computed from rows up through $t$ only.
The uncompressed model's logits and every oracle divergence field are
excluded from the input space by construction.
\todo{Cite the frozen feature-inventory hash from the protocol lock.}

\paragraph{Rate-based multi-horizon architecture.}
Predicting the four cumulative targets independently could produce
incoherent forecasts, such as less predicted damage at 50 tokens than at 25.
HERALD instead predicts four nonnegative average damage \emph{rates} over the
disjoint segments $1{:}5$, $6{:}10$, $11{:}25$, and $26{:}50$ of the forecast
window, one XGBoost regressor per segment, and reconstructs the cumulative
forecasts as
\begin{equation}
\hat Y_5 = 5R_1, \quad
\hat Y_{10} = \hat Y_5 + 5R_2, \quad
\hat Y_{25} = \hat Y_{10} + 15R_3, \quad
\hat Y_{50} = \hat Y_{25} + 25R_4.
\label{eq:rates}
\end{equation}
Because the rates are clamped nonnegative, the horizon forecasts are monotone
by construction. When no compression is active, the target is identically
zero by definition and the system returns zero without invoking the learned
model.

\paragraph{Model selection.}
A causal temporal convolutional network was trained as the sequence-model
alternative under the same protocol. The selection rule, frozen before the
confirmation set was opened, chose by development macro MSE: XGBoost at
\todo{0.003736; verify against protocol lock} against the TCN at
\todo{0.003801}, with the TCN worse at every horizon. The snapshot model won;
sequence modeling added no measurable value over causally aggregated history
features at this scale.

\paragraph{Comparators.}
The primary comparator is an action-clock predictor: the strongest
non-causal reference that knows the compression action and the generation
position but nothing about the running distribution.
\todo{State its exact construction from the protocol lock.} Improvement over
this comparator isolates the value of the online logit signal, since action
and position are available to any scheduler without observing the stream.
